% Source: https://artofproblemsolving.com/community/c3h2692438_sequences
% Stable identity: aops-topic-2692438-q01
For $n\ge1$ let $f(n)=\sqrt{p^2n^2+1}-pn=\dfrac{1}{\sqrt{p^2n^2+1}+pn}$ (this is $\sqrt m-pn$ with $m=p^2n^2+1$). Then $f$ is positive, strictly decreasing and tends to $0$. Moreover $f(n)<2f(n+1)$, because
\[\sqrt{p^2(n+1)^2+1}+p(n+1)<2\sqrt{p^2n^2+1}+2pn:\]
indeed $p(n+1)\le2pn$, and $4(p^2n^2+1)-\left(p^2(n+1)^2+1\right)=p^2(n-1)(3n+1)+3>0$.

\emph{Case 1: $2\varepsilon\le f(1)$.} Let $n$ be the least positive integer with $f(n)<2\varepsilon$. Then $n\ge2$ and $f(n-1)\ge2\varepsilon$, so $f(n)>\dfrac{f(n-1)}{2}\ge\varepsilon$. Thus $m=p^2n^2+1$ works.

\emph{Case 2: $f(1)<2\varepsilon$.} Take $n=1$ and let $h(k)=\sqrt{p^2+k}-p$ for $k\ge1$; $h$ increases without bound and $h(1)=f(1)$. Let $k$ be the least index with $h(k)>\varepsilon$. If $k=1$, then $\varepsilon<f(1)<2\varepsilon$ and $m=p^2+1$ works. If $k\ge2$, then $h(k-1)\le\varepsilon$, and since the gaps $\sqrt{N+1}-\sqrt N=\dfrac{1}{\sqrt{N+1}+\sqrt N}$ decrease as $N$ grows,
\[h(k)-h(k-1)=\sqrt{p^2+k}-\sqrt{p^2+k-1}<\sqrt{p^2+1}-p=h(1)\le h(k-1)\le\varepsilon.\]
Hence $\varepsilon<h(k)<2\varepsilon$, and $m=p^2+k$, $n=1$ work.

In every case suitable positive integers $m,n$ exist. $\blacksquare$
